\chapter{Random Sampling Algorithms on Data Streams}
\label{chap:random-sampling-data-streams}

\begin{center}
  \textbf{Lecture 7}\\
  \emph{Random Sampling Algorithms on Data Streams: Streaming Model, Uniform Sampling, Algorithm S, Priority Sampling, and Reservoir Sampling}\\[0.5em]
  Date: October 6, 2026
\end{center}

In previous lectures, we thoroughly examined data processing under the \textit{External Memory Model}, focusing on massive datasets stored on secondary storage through multiway external merge sort and distribution sort algorithms. In those scenarios, although access costs were dominated by block I/O transfers, multiple scans (\textit{multi-pass}) over the stored dataset were still permissible.

In this chapter, we explore a fundamental paradigm shift: the \textbf{Data Streaming Model}. In this setting, data arrives as a continuous, massive, or potentially infinite sequence at high arrival rates, subject to stringent working memory and per-item processing constraints. Among the fundamental stream processing tasks, \textbf{Uniform Random Sampling} plays a critical role in constructing statistical synopses, network monitoring, high-volume log debugging, and online aggregate estimation.

We investigate both principal categories of sampling without replacement: when stream length $n$ is known in advance (analyzing single-item selection and Knuth's \textbf{Algorithm S}, supported by a complete combinatorial derivation) and when cardinality $n$ is unknown ahead of time (\textbf{Priority Sampling} with Min-Heaps and the Waterman--Knuth \textbf{Reservoir Sampling} framework). We prove the correctness of Reservoir Sampling by mathematical induction, present comprehensive step-by-step traces, and examine engineering optimizations (including Vitter's \textbf{Algorithm Z}).

\FloatBarrier
\section{Foundations of the Streaming Model and Random Sampling}
\label{sec:streaming-model-foundations}

\subsection{The Data Streaming Computation Model}
In traditional computational paradigms (both in RAM and secondary storage), datasets remain available for repeated scans. Under the \textit{Data Streaming Model}, conversely, the algorithm can neither control data ordering nor request data retransmission.

The streaming scenario is formally characterized by four architectural pillars:
\begin{enumerate}
  \item \textbf{Data Stream $S$:} a time-ordered sequence of discrete items, potentially infinite or of massive size $n$ ($n \to \infty$ or $n \gg M$):
  \begin{equation}
    S = \langle i_1, i_2, \dots, i_j, \dots, i_n \rangle.
  \end{equation}
  \item \textbf{Strict Working Memory Constraint ($M \ll n$):} the fast internal working memory $M$ available to the algorithm satisfies $M \ll n$. Typically, storage is strictly bounded by:
  \begin{equation}
    M = \mathcal{O}(m) \quad \text{or} \quad M = \mathcal{O}(\operatorname{polylog}(n)) \text{ words of memory},
  \end{equation}
  where $m$ denotes the desired sample size. Consequently, storing the entire stream or any substantial fraction thereof is impossible.
  \item \textbf{Single-Pass Sequential Access:} elements arrive one by one in chronological order. Once an item $i_j$ has passed, if the algorithm chooses not to store it in working memory, it is \textbf{lost forever} and can never be re-examined.
  \item \textbf{Strictly Bounded Per-Item Processing Time:} to avoid input buffer exhaustion and packet drops on high-speed feeds (such as backbone routers or real-time trading systems), the processing time allocated to each arriving item $i_j$ must remain strictly constant:
  \begin{equation}
    T(i_j) = \mathcal{O}(1) \quad \text{or at most} \quad T(i_j) = \mathcal{O}(\log m).
  \end{equation}
\end{enumerate}

\subsection{Formal Definition of the Uniform Random Sampling Problem}
Given a continuous stream $S$, our objective is to select and maintain a subset of $m$ items that represents a \textbf{Simple Random Sample without Replacement} (SRSWOR).

Formally, at any step $n$ (or upon stream completion if $n$ is finite), two probabilistic uniformity criteria must hold:

\begin{definition}[Marginal Uniformity]
Every single element $i_j \in S$ (for $1 \le j \le n$) must possess an identical marginal inclusion probability:
\begin{equation}
  \mathbb{P}(i_j \in \text{Sample}) = \frac{m}{n}.
  \label{eq:marginal-uniformity}
\end{equation}
\end{definition}

\begin{definition}[Joint Uniformity]
Every possible subset $S^* \subseteq S$ of cardinality $m$ must have an identical joint extraction probability:
\begin{equation}
  \mathbb{P}(\text{Sample} = S^*) = \frac{1}{\binom{n}{m}}, \qquad \forall S^* \subseteq S \text{ such that } |S^*| = m.
  \label{eq:joint-uniformity}
\end{equation}
\end{definition}

Depending on prior knowledge of stream cardinality $n$, sampling algorithms are classified into two primary regimes:
\begin{itemize}
  \item \textbf{Case A ($n$ known a priori):} the total stream length $n$ is known before sampling begins (e.g., sequentially reading an unindexed $n$-line log file from tape or storage).
  \item \textbf{Case B ($n$ unknown a priori):} cardinality $n$ is unavailable upfront, and the stream may terminate arbitrarily or continue indefinitely. The algorithm must guarantee that whenever the stream terminates at any step $t \ge m$, internal memory stores a valid uniform random sample drawn from the first $t$ items.
\end{itemize}

\FloatBarrier
\section{Sampling with Known Stream Length (\texorpdfstring{$n$}{n} Known)}
\label{sec:sampling-known-n}

\subsection{Single-Item Selection (\texorpdfstring{$m = 1$}{m = 1})}
Consider the elementary case of extracting exactly one element ($m = 1$) from a stream of known length $n$, ensuring each item has uniform probability $\frac{1}{n}$ of being selected while deciding online whether to accept or reject the current item $i_j$.

\subsubsection{Algorithm Formulation}
The algorithm sequentially inspects elements for $j = 1, 2, \dots, n$:
\begin{enumerate}
  \item At step $j$, draw a random variable $p \sim \operatorname{Uniform}(0, 1)$, meaning $p \in [0, 1)$.
  \item Evaluate the selection condition:
  \begin{equation}
    \text{If } p \le \frac{1}{n - j + 1}, \quad \text{select } i_j \text{ and terminate the algorithm.}
    \label{eq:single-item-threshold}
  \end{equation}
  \item If $p > \frac{1}{n - j + 1}$, item $i_j$ is permanently discarded, and processing proceeds to step $j+1$.
\end{enumerate}

\begin{algorithm}[H]
  \caption{Online Single-Item Sampling ($m = 1$, $n$ known)}
  \label{alg:single-element-sampling}
  \KwInput{Stream $S = \langle i_1, \dots, i_n \rangle$ of known length $n$}
  \KwOutput{A single item $x \in S$ sampled uniformly with $\mathbb{P}(x = i_j) = \frac{1}{n}$}
  \For{$j \leftarrow 1$ \KwTo $n$}{
    $p \leftarrow \operatorname{Uniform}(0, 1)$\;
    \If{$p \le \frac{1}{n - j + 1}$}{
      \Return $i_j$\;
    }
  }
\end{algorithm}

\subsubsection{Rigorous Proof by Induction}
\begin{theorem}
Algorithm~\ref{alg:single-element-sampling} ensures that every element $i_j \in S$ has exact probability $\frac{1}{n}$ of being chosen:
\begin{equation}
  \mathbb{P}(\text{select } i_j) = \frac{1}{n}, \qquad \forall j \in \{1, 2, \dots, n\}.
\end{equation}
\end{theorem}

\begin{proof}
We proceed by induction on stream index $j$.

\textbf{Base step ($j = 1$):}\\
At step 1, for item $i_1$, the selection threshold is:
\begin{equation}
  \frac{1}{n - 1 + 1} = \frac{1}{n}.
\end{equation}
Since $p \sim \operatorname{Uniform}(0, 1)$, the probability that $p \le \frac{1}{n}$ matches the interval length:
\begin{equation}
  \mathbb{P}(\text{pick } i_1) = \frac{1}{n}.
\end{equation}
The base case holds.

\textbf{Inductive hypothesis:}\\
Assume the property holds for all preceding indices $k \in \{1, 2, \dots, j - 1\}$, meaning:
\begin{equation}
  \mathbb{P}(\text{pick } i_k) = \frac{1}{n}.
\end{equation}
Because the algorithm halts immediately upon selecting an element, the events $\{\text{pick } i_k\}$ are mutually disjoint. Thus, the cumulative probability that at least one of the first $j - 1$ elements was selected is:
\begin{equation}
  \mathbb{P}\left(\bigcup_{k=1}^{j-1} \{\text{pick } i_k\}\right) = \sum_{k=1}^{j-1} \mathbb{P}(\text{pick } i_k) = \sum_{k=1}^{j-1} \frac{1}{n} = \frac{j - 1}{n}.
\end{equation}
The complementary probability that \emph{no} element among the first $j - 1$ was selected is:
\begin{equation}
  \mathbb{P}(\text{no element selected in } \{i_1, \dots, i_{j-1}\}) = 1 - \frac{j - 1}{n} = \frac{n - j + 1}{n}.
  \label{eq:prob-none-selected-before}
\end{equation}

\textbf{Inductive step:}\\
For element $i_j$ to be chosen at step $j$, two joint events must occur:
\begin{enumerate}
  \item No element among $i_1, \dots, i_{j-1}$ was selected earlier.
  \item At step $j$, the random draw $p$ satisfies $p \le \frac{1}{n - j + 1}$.
\end{enumerate}
\begin{align}
  \mathbb{P}(\text{pick } i_j) &= \mathbb{P}\left(p \le \frac{1}{n - j + 1} \;\middle|\; \text{no element chosen previously}\right) \notag\\
  &\quad \cdot \mathbb{P}(\text{no element chosen previously}) \notag\\
  &= \frac{1}{n - j + 1} \cdot \left(\frac{n - j + 1}{n}\right) = \frac{1}{n}.
\end{align}
The factor $(n - j + 1)$ cancels exactly, yielding $\frac{1}{n}$. The assertion holds for all $j \in \{1, \dots, n\}$.
\end{proof}

\subsubsection{Geometric Interpretation}
The probabilistic mechanism underlying Algorithm~\ref{alg:single-element-sampling} admits a clear geometric interpretation illustrated in \autoref{fig:streaming-partition-m1}.

\begin{figure}[H]
  \centering
  \resizebox{0.95\textwidth}{!}{%
  \begin{tikzpicture}[
    scale=1.0,
    box/.style={rectangle, draw=blue!80!black, thick, minimum height=1.3cm, align=center, font=\small},
    brace/.style={decorate, decoration={brace, amplitude=6pt, raise=4pt}, thick}
  ]
    % Segment processed elements
    \node[box, fill=gray!15, minimum width=4.8cm, text width=4.6cm] (past) at (0, 0) {
      \textbf{Steps $1 \dots j - 1$ ($j - 1$ items)}\\[2pt]
      \footnotesize Evaluated: none selected
    };

    % Current element
    \node[box, fill=teal!20, draw=teal!80!black, minimum width=2.8cm, text width=2.6cm, right=0.35cm of past] (curr) {
      \textbf{Step $j$}: item $i_j$\\[2pt]
      \footnotesize $\mathbb{P} = \frac{1}{n-j+1}$
    };

    % Future residual elements
    \node[box, fill=blue!10, minimum width=5.6cm, text width=5.4cm, right=0.35cm of curr] (future) {
      \textbf{Steps $j + 1 \dots n$ ($n - j$ items)}\\[2pt]
      \footnotesize Remaining future items awaiting inspection
    };

    % Brace for remaining pool above boxes
    \draw[brace, draw=teal!80!black] ($(curr.north west)+(0,0.1)$) -- ($(future.north east)+(0,0.1)$)
      node[midway, above=8pt, font=\small\bfseries, text=teal!80!black] {
        Residual candidate pool: $n - j + 1$ equally likely elements
      };

    % Brace total stream below boxes
    \draw[brace, decoration={mirror, amplitude=6pt, raise=4pt}, draw=blue!85!black] ($(past.south west)+(0,-0.15)$) -- ($(future.south east)+(0,-0.15)$)
      node[midway, below=12pt, font=\small\bfseries, text=blue!85!black] {
        Full Stream of known cardinality $n$
      };
  \end{tikzpicture}%
  }
  \caption{Stream space partition at step $j$ for single-element sampling ($m = 1$).}
  \label{fig:streaming-partition-m1}
\end{figure}

If no element was selected across the first $j - 1$ steps, the target item must reside within the remaining $n - j + 1$ elements. By uniformity, all remaining items share equal rights to selection, so the conditional probability of choosing $i_j$ immediately is the reciprocal of the remaining count: $\frac{1}{n - j + 1}$.

\FloatBarrier
\subsection{Multi-Item Selection: Algorithm S (Fan, Muller, Rezucha / Knuth)}
When the requested sample size is $m > 1$ and total stream size $n$ is known upfront, this principle generalizes to \textbf{Algorithm S} (introduced by Fan, Muller, and Rezucha~\cite{fan1962} and popularized by Knuth~\cite{knuth1997}), which maintains a dynamic counter of selected elements.

\subsubsection{State Variable Definitions}
The algorithm tracks the following quantities during sequential traversal:
\begin{itemize}
  \item $n$: total number of stream elements;
  \item $m$: target sample size;
  \item $j \in \{1, 2, \dots, n\}$: time index of current candidate element $i_j$;
  \item $s$: number of items already chosen into the sample across the first $j - 1$ steps ($0 \le s \le m$);
  \item $m - s$: remaining items \emph{still needed} to complete the sample;
  \item $n - j + 1$: total remaining items in the stream yet to be evaluated (including $i_j$).
\end{itemize}

\subsubsection{Selection Rule and Boundary Cases}
At each iteration $j \in \{1, \dots, n\}$:
\begin{enumerate}
  \item Draw $p \sim \operatorname{Uniform}(0, 1)$.
  \item Compare $p$ against the dynamic threshold $\frac{m - s}{n - j + 1}$:
  \begin{equation}
    p \le \frac{m - s}{n - j + 1}.
    \label{eq:algorithm-s-threshold}
  \end{equation}
  \item If satisfied: insert $i_j$ into sample and increment counter:
  \begin{equation}
    s \leftarrow s + 1.
  \end{equation}
  Otherwise, discard $i_j$ and keep $s$ unchanged.
\end{enumerate}

The algorithm handles boundary conditions automatically:
\begin{itemize}
  \item \textbf{Sample Completed ($s = m$):} the threshold numerator becomes $m - s = 0$. Selection probability falls to zero ($\frac{0}{n - j + 1} = 0$), preventing further additions and enabling early termination.
  \item \textbf{Forced Inclusion of All Remaining Items ($m - s = n - j + 1$):} the remaining quota matches the remaining elements. The threshold becomes $\frac{m - s}{n - j + 1} = 1$, deterministically forcing inclusion of all remaining elements since $p < 1$ with probability $1$.
\end{itemize}

\begin{algorithm}[H]
  \caption{Algorithm S (Selection Sampling by Fan, Muller, Rezucha / Knuth)}
  \label{alg:algorithm-s}
  \KwInput{Stream $S = \langle i_1, \dots, i_n \rangle$ of known length $n$, sample size $m \le n$}
  \KwOutput{Set $C$ containing $m$ items sampled uniformly}
  $C \leftarrow \emptyset$\;
  $s \leftarrow 0$\tcp*{Counter of already sampled items}
  \For{$j \leftarrow 1$ \KwTo $n$}{
    \If{$s = m$}{
      \textbf{break}\tcp*{Sample full: terminate early}
    }
    $p \leftarrow \operatorname{Uniform}(0, 1)$\;
    \If{$p \le \frac{m - s}{n - j + 1}$}{
      $C \leftarrow C \cup \{i_j\}$\;
      $s \leftarrow s + 1$\;
    }
  }
  \Return $C$\;
\end{algorithm}

\subsubsection{Combinatorial Derivation of the Threshold}
The threshold $\frac{m - s}{n - j + 1}$ stems directly from combinatorial principles and the hypergeometric distribution.

\begin{theorem}
Given that $s$ elements have already been selected across the first $j - 1$ steps, the conditional probability that the current item $i_j$ belongs to a uniform sample of size $m$ drawn from the remaining $n - j + 1$ elements is:
\begin{equation}
  \mathbb{P}(i_j \in \text{Sample} \mid s \text{ already sampled}) = \frac{m - s}{n - j + 1}.
\end{equation}
\end{theorem}

\begin{proof}
The conditional probability is the ratio of favorable combinations to total configurations:
\begin{equation}
  \mathbb{P}(\text{pick } i_j \mid s) = \frac{\text{Favorable cases (complete sample with } i_j)}{\text{Possible cases (choose } m - s \text{ from residual elements)}}.
\end{equation}

Fixing $i_j$ accounts for $1$ item; thus, $m - s - 1$ items must be chosen from the subsequent $n - j$ available stream items. The number of ways is given by the binomial coefficient:
\begin{equation}
  \binom{n - j}{m - s - 1}.
\end{equation}
Conversely, the unconstrained combinations for choosing $m - s$ items from all $n - j + 1$ remaining candidates is:
\begin{equation}
  \binom{n - j + 1}{m - s}.
\end{equation}

Computing the ratio and expanding factorials:
\begin{align}
  \frac{\binom{n - j}{m - s - 1}}{\binom{n - j + 1}{m - s}} &= \frac{\dfrac{(n - j)!}{(m - s - 1)! \, ((n - j) - (m - s - 1))!}}{\dfrac{(n - j + 1)!}{(m - s)! \, ((n - j + 1) - (m - s))!}} \notag\\[1em]
  &= \frac{\dfrac{(n - j)!}{(m - s - 1)! \, (n - j - m + s + 1)!}}{\dfrac{(n - j + 1)!}{(m - s)! \, (n - j - m + s + 1)!}}.
\end{align}
The factor $(n - j - m + s + 1)!$ in the denominator of both fractions cancels completely. Regrouping:
\begin{equation}
  = \frac{(n - j)!}{(n - j + 1)!} \cdot \frac{(m - s)!}{(m - s - 1)!}.
\end{equation}
Applying the recursive factorial definition $(k! = k \cdot (k-1)!)$:
\begin{align}
  (n - j + 1)! &= (n - j + 1) \cdot (n - j)!, \\
  (m - s)! &= (m - s) \cdot (m - s - 1)!.
\end{align}
Substituting into each term:
\begin{equation}
  = \frac{1}{n - j + 1} \cdot (m - s) = \frac{m - s}{n - j + 1}.
\end{equation}
Thus, the local transition probability exactly matches the global hypergeometric uniform distribution.
\end{proof}

\FloatBarrier
\subsection{Step-by-Step Walkthrough: Algorithm S Example}
To examine how stochastic thresholds operate in practice, we trace the example from lecture notes.

\begin{example}[Algorithm S Execution Trace]
Consider the following input configuration:
\begin{itemize}
  \item Input stream: $S = \langle a, b, c, d, e, f, g, h \rangle \implies n = 8$.
  \item Target sample size: $m = 2$.
  \item Sequence of pseudo-random draws $p \in [0, 1)$ drawn per step:
  \begin{equation}
    p = \langle 0.5, \; 0.5, \; 0.1, \; 0.5, \; 0.5, \; 0.2, \; 0.9, \; 0.3 \rangle.
  \end{equation}
\end{itemize}

The step-by-step state progression is presented in \autoref{tab:trace-algorithm-s}.

\begin{table}[H]
  \centering
  \caption{Step-by-step execution trace of Algorithm S ($n = 8$, $m = 2$).}
  \label{tab:trace-algorithm-s}
  \resizebox{\textwidth}{!}{%
  \begin{tabular}{c c c c c c c c c c}
    \toprule
    $j$ & $i_j$ & $s$ & $m - s$ & $n - j + 1$ & Threshold $\frac{m-s}{n-j+1}$ & $p$ & Condition $p \le \text{Threshold}$ & Outcome & New $s$ \\
    \midrule
    $1$ & $a$ & $0$ & $2$ & $8$ & $\frac{2}{8} = 0.250$ & $0.5$ & $0.5 \le 0.250$ (False) & Discarded ($\times$) & $0$ \\
    $2$ & $b$ & $0$ & $2$ & $7$ & $\frac{2}{7} \approx 0.286$ & $0.5$ & $0.5 \le 0.286$ (False) & Discarded ($\times$) & $0$ \\
    $3$ & $c$ & $0$ & $2$ & $6$ & $\frac{2}{6} \approx 0.333$ & $0.1$ & $0.1 \le 0.333$ (True) & \textbf{Selected ($\checkmark$)} & $1$ \\
    $4$ & $d$ & $1$ & $1$ & $5$ & $\frac{1}{5} = 0.200$ & $0.5$ & $0.5 \le 0.200$ (False) & Discarded ($\times$) & $1$ \\
    $5$ & $e$ & $1$ & $1$ & $4$ & $\frac{1}{4} = 0.250$ & $0.5$ & $0.5 \le 0.250$ (False) & Discarded ($\times$) & $1$ \\
    $6$ & $f$ & $1$ & $1$ & $3$ & $\frac{1}{3} \approx 0.333$ & $0.2$ & $0.2 \le 0.333$ (True) & \textbf{Selected ($\checkmark$)} & $2$ \\
    $7$ & $g$ & $2$ & $0$ & $2$ & $\frac{0}{2} = 0.000$ & $0.9$ & $0.9 \le 0.000$ (False) & Discarded / STOP & $2$ \\
    $8$ & $h$ & $2$ & $0$ & $1$ & $\frac{0}{1} = 0.000$ & $0.3$ & $0.3 \le 0.000$ (False) & Discarded / STOP & $2$ \\
    \bottomrule
  \end{tabular}%
  }
\end{table}

\textbf{Outcome:} At step $j = 6$, the counter reaches $s = m = 2$, completing the sample $C = \{c, f\}$. At steps $7$ and $8$, the selection threshold drops to zero, systematically discarding subsequent elements immediately.
\end{example}

\FloatBarrier
\section{Sampling with Unknown Stream Length (\texorpdfstring{$n$}{n} Unknown)}
\label{sec:sampling-unknown-n}

In realistic streaming environments—including network telemetry, server logging, and sensor feeds—the stream size $n$ is unknown upfront and data may cease arriving at any arbitrary time. Streaming algorithms must therefore maintain a valid sample at all moments.

\subsection{Priority Sampling with Min-Heap (Random Keys Sampling)}
A natural strategy for sampling without knowing $n$ relies on \textbf{Random Keys Sampling} (or \textit{Priority Sampling}).

\subsubsection{Functional Principle}
Each arriving item $i_j$ receives an independent and identically distributed continuous random key:
\begin{equation}
  \pi_j \sim \operatorname{Uniform}(0, 1).
\end{equation}
The algorithm dynamically tracks the $m$ elements with the highest keys observed so far. Because keys are continuous and i.i.d., every subset of size $m$ from the first $j$ items has equal likelihood of possessing the $m$ largest keys.

\subsubsection{Management via Bounded Min-Heap}
To confine memory consumption to $\mathcal{O}(m)$, we maintain a \textbf{Min-Heap} $H$ of fixed capacity $m$, storing pairs $\langle \pi_j, i_j \rangle$ keyed by $\pi_j$:
\begin{itemize}
  \item \textbf{First $m$ elements ($j \le m$):} insert pairs directly into $H$. Total initial construction time is $\mathcal{O}(m)$.
  \item \textbf{Subsequent elements ($j > m$):} for each candidate $i_j$ with key $\pi_j$:
  \begin{enumerate}
    \item Inspect the heap root, which holds the current minimum priority among the $m$ sample elements:
    \begin{equation}
      \pi_{\min} = \operatorname{min\_key}(H).
    \end{equation}
    \item \textbf{If $\pi_j \le \pi_{\min}$:} the key fails to enter the top $m$. Item $i_j$ is discarded immediately in $\mathcal{O}(1)$ time.
    \item \textbf{If $\pi_j > \pi_{\min}$:} the element holding $\pi_{\min}$ is ejected. We replace the root (\texttt{Replace-Min}) with $\langle \pi_j, i_j \rangle$ and restore heap order via \texttt{Heapify-Down} in $\mathcal{O}(\log m)$ time.
  \end{enumerate}
\end{itemize}

The functional architecture of Priority Sampling is illustrated in \autoref{fig:priority-sampling-minheap}.

\begin{figure}[H]
  \centering
  \resizebox{0.95\textwidth}{!}{%
  \begin{tikzpicture}[
    scale=0.95,
    box/.style={rectangle, draw=blue!85!black, fill=blue!10, thick, minimum width=2.4cm, minimum height=1.1cm, align=center, font=\small},
    heapnode/.style={circle, draw=orange!90!black, fill=orange!15, thick, minimum size=1.0cm, align=center, font=\scriptsize\bfseries},
    rootnode/.style={circle, draw=red!85!black, fill=red!20, line width=1.5pt, minimum size=1.1cm, align=center, font=\scriptsize\bfseries},
    decision/.style={diamond, draw=teal!85!black, fill=teal!15, thick, aspect=1.8, align=center, font=\footnotesize\bfseries, inner sep=2pt},
    arrow/.style={-{Stealth[scale=0.95]}, thick, draw=blue!80!black}
  ]
    % Stream
    \node[box, fill=gray!15, draw=gray!70!black] (stream) at (0, 0) {
      Stream $S$\\
      \footnotesize Item $i_j$
    };

    % Key Generator
    \node[box, fill=teal!15, draw=teal!80!black, right=1.0cm of stream] (rng) {
      Key Generator\\
      \footnotesize $\pi_j \sim \mathcal{U}(0, 1)$
    };

    % Decision
    \node[decision, right=1.2cm of rng] (dec) {
      $\pi_j > \pi_{\min}$?
    };

    % Drop
    \node[box, fill=red!10, draw=red!70!black, below=1.4cm of dec] (drop) {
      Discard Item\\
      \footnotesize Cost $\mathcal{O}(1)$
    };

    % Min-Heap Container
    \node[draw=orange!80!black, dashed, rounded corners, fill=orange!5, minimum width=4.8cm, minimum height=3.6cm, right=2.0cm of dec] (heapbox) {};
    \node[above=3pt of heapbox.north, font=\small\bfseries, text=orange!90!black] {Min-Heap $H$ (Capacity $m$)};

    % Heap nodes
    \node[rootnode] (r) at ($(heapbox.north)+(0, -1.0)$) {$\pi_{\min}$\\\tiny (Root)};
    \node[heapnode] (c1) at ($(r)+(-1.4, -1.2)$) {$\pi_{a}$};
    \node[heapnode] (c2) at ($(r)+(1.4, -1.2)$) {$\pi_{b}$};

    \draw[thick, draw=orange!80!black] (r) -- (c1);
    \draw[thick, draw=orange!80!black] (r) -- (c2);

    % Flow arrows
    \draw[arrow] (stream) -- (rng);
    \draw[arrow] (rng) -- node[above=1pt, font=\scriptsize] {$\langle \pi_j, i_j \rangle$} (dec);
    \draw[arrow] (dec) -- node[right=1pt, font=\scriptsize\bfseries, text=red!80!black] {NO} (drop);
    \draw[arrow] (dec.east) -- ++(0.9, 0) |- node[pos=0.15, above=2pt, font=\scriptsize\bfseries, text=teal!80!black] {YES} node[pos=0.82, below=2pt, font=\scriptsize] {\texttt{Replace-Min}} (r.west);
  \end{tikzpicture}%
  }
  \caption{Functional flow of streaming Priority Sampling using a bounded Min-Heap of capacity $m$.}
  \label{fig:priority-sampling-minheap}
\end{figure}

\begin{property}[Complexity Analysis of Priority Sampling]
Memory space is bounded by $\mathcal{O}(m)$ words. Worst-case processing time per element is $\mathcal{O}(\log m)$. However, because the probability that element $j$ qualifies into the top $m$ is $\frac{m}{j}$, the expected number of heap updates across $n$ stream items is:
\begin{equation}
  \sum_{j=m+1}^n \frac{m}{j} = m \sum_{j=m+1}^n \frac{1}{j} = \mathcal{O}(m \log(n/m)).
\end{equation}
The expected amortized cost per item is:
\begin{equation}
  T_{\text{avg}} = \mathcal{O}\!\left(1 + \frac{m \log m \log(n/m)}{n}\right) = \mathcal{O}(1) \quad \text{for } n \gg m.
\end{equation}
Nevertheless, Priority Sampling requires generating a pseudo-random key for \emph{every} stream item, alongside tree management overhead.
\end{property}

\FloatBarrier
\subsection{Reservoir Sampling (Algorithm R by Waterman / Knuth)}
To achieve strict $\mathcal{O}(1)$ worst-case time per item without priority heap overhead, the canonical solution is \textbf{Reservoir Sampling} (originally developed by Alan G. Waterman and formalized as \textbf{Algorithm R} by Knuth~\cite{knuth1997}).

\subsubsection{Operational Mechanism}
The algorithm maintains a simple linear array $R[1 \dots m]$, called the \textbf{reservoir}. Execution proceeds in two phases:
\begin{enumerate}
  \item \textbf{Phase 1 (Deterministic Initialization):} the first $m$ stream items are unconditionally assigned to their matching slots in $R$:
  \begin{equation}
    R[j] \leftarrow i_j, \qquad \forall j \in \{1, 2, \dots, m\}.
  \end{equation}
  No pseudo-random numbers are generated during this initial phase.
  \item \textbf{Phase 2 (Competitive Online Replacement):} for every subsequent item $i_j$ arriving at time $j = m + 1, m + 2, \dots$:
  \begin{enumerate}
    \item Draw a uniform discrete random integer $k$ within the inclusive range $[1, j]$:
    \begin{equation}
      k \sim \operatorname{rand}(1, j).
    \end{equation}
    \item Evaluate admission:
    \begin{itemize}
      \item If $k \le m$: item $i_j$ enters the reservoir, \textbf{replacing} slot $k$:
      \begin{equation}
        R[k] \leftarrow i_j.
      \end{equation}
      The item formerly residing at $R[k]$ is permanently evicted.
      \item If $k > m$: current item $i_j$ is discarded immediately, leaving $R$ unchanged.
    \end{itemize}
  \end{enumerate}
\end{enumerate}

\begin{algorithm}[H]
  \caption{Reservoir Sampling (Algorithm R by Waterman / Knuth)}
  \label{alg:reservoir-sampling-stream}
  \KwInput{Streaming sequence $S = \langle i_1, i_2, \dots \rangle$, reservoir capacity $m$}
  \KwOutput{Array $R[1 \dots m]$ containing a uniform sample at all times}
  \For{$j \leftarrow 1$ \KwTo $m$}{
    $R[j] \leftarrow i_j$\tcp*{Fill initial reservoir without RNG}
  }
  $j \leftarrow m + 1$\;
  \While{Stream continues and new item $i_j$ arrives}{
    $k \leftarrow \operatorname{RandomInteger}(1, j)$\;
    \If{$k \le m$}{
      $R[k] \leftarrow i_j$\tcp*{Replace slot k with probability m/j}
    }
    $j \leftarrow j + 1$\;
  }
  \Return $R$\;
\end{algorithm}

The replacement mechanism is illustrated in \autoref{fig:reservoir-sampling-flow}.

\begin{figure}[H]
  \centering
  \resizebox{0.95\textwidth}{!}{%
  \begin{tikzpicture}[
    scale=0.95,
    box/.style={rectangle, draw=blue!80!black, fill=blue!10, thick, minimum width=2.4cm, minimum height=1.1cm, align=center, font=\small},
    slot/.style={rectangle, draw=blue!85!black, thick, minimum width=1.1cm, minimum height=1.1cm, align=center, font=\small\bfseries},
    replaced/.style={rectangle, draw=teal!85!black, fill=teal!25, line width=1.5pt, minimum width=1.1cm, minimum height=1.1cm, align=center, font=\small\bfseries},
    decision/.style={diamond, draw=teal!85!black, fill=teal!15, thick, aspect=1.8, align=center, font=\footnotesize\bfseries, inner sep=2pt},
    arrow/.style={-{Stealth[scale=0.95]}, thick, draw=blue!80!black}
  ]
    % Stream
    \node[box, fill=gray!15, draw=gray!70!black] (stream) at (0, 0) {
      Stream $S$\\
      \footnotesize Item $i_j$ ($j > m$)
    };

    % Integer RNG
    \node[box, fill=teal!15, draw=teal!80!black, right=1.0cm of stream] (rng) {
      RNG $k$\\
      \footnotesize $k \in \{1, \dots, j\}$
    };

    % Decision
    \node[decision, right=1.1cm of rng] (dec) {
      $k \le m$?
    };

    % Discard
    \node[box, fill=red!10, draw=red!70!black, below=1.3cm of dec] (drop) {
      Discard $i_j$\\
      \footnotesize Probability $1 - \frac{m}{j}$
    };

    % Reservoir Array
    \node[slot, fill=blue!8, right=2.2cm of dec] (s1) {$R[1]$};
    \node[slot, fill=blue!8, right=0pt of s1] (s2) {$R[2]$};
    \node[slot, fill=gray!15, right=0pt of s2] (dots) {$\dots$};
    \node[replaced, right=0pt of dots] (sk) {$i_j$};
    \node[slot, fill=blue!8, right=0pt of sk] (sm) {$R[m]$};

    \node[below=8pt of dots, font=\small\bfseries, text=blue!85!black] {Reservoir Memory Array $R[1 \dots m]$};

    % Arrows
    \draw[arrow] (stream) -- (rng);
    \draw[arrow] (rng) -- node[above=1pt, font=\scriptsize] {$k$} (dec);
    \draw[arrow] (dec) -- node[right=1pt, font=\scriptsize\bfseries, text=red!80!black] {NO} (drop);
    \draw[arrow, draw=teal!85!black, very thick] (dec.north) -- ++(0, 0.8) -| node[pos=0.25, above=2pt, font=\scriptsize\bfseries, text=teal!80!black] {YES (Probability $\frac{m}{j}$)} (sk.north)
      node[pos=0.88, right=3pt, font=\footnotesize\bfseries, text=teal!90!black] {Slot $R[k] \leftarrow i_j$};
  \end{tikzpicture}%
  }
  \caption{Operational replacement flow of Reservoir Sampling (Algorithm R).}
  \label{fig:reservoir-sampling-flow}
\end{figure}

\FloatBarrier
\subsection{Step-by-Step Walkthrough: Reservoir Sampling Example}
We trace the step-by-step state progression on a sample stream.

\begin{example}[Reservoir Sampling Execution Trace]
Consider the sequence:
\begin{equation}
  S = \langle a, b, c, d, e, f, g, h, i \rangle, \qquad \text{with capacity } m = 3.
\end{equation}
\begin{itemize}
  \item \textbf{Initialization ($j = 1, 2, 3$):} the first $m = 3$ elements are saved into $R$:
  \begin{equation}
    R = [a, \; b, \; c].
  \end{equation}
  \item \textbf{Online updates ($j = 4, 5, 6, 7$):} state transitions are detailed in \autoref{tab:trace-reservoir-sampling}.
\end{itemize}

\begin{table}[H]
  \centering
  \caption{Step-by-step trace of Reservoir Sampling ($m = 3$).}
  \label{tab:trace-reservoir-sampling}
  \resizebox{\textwidth}{!}{%
  \begin{tabular}{c c c c l c}
    \toprule
    Step $j$ & Item $i_j$ & Random Integer $k$ & $k \le m$ & Executed Action & Reservoir State $R[1 \dots 3]$ \\
    \midrule
    $1 \dots 3$ & $a, b, c$ & --- & --- & Deterministic initialization & $[a, b, c]$ \\
    \addlinespace
    $4$ & $d$ & $k = 2 \in \{1 \dots 4\}$ & True ($2 \le 3$) & Replace slot $R[2] \leftarrow d$ & $[a, d, c]$ \\
    \addlinespace
    $5$ & $e$ & $k = 5 \in \{1 \dots 5\}$ & False ($5 > 3$) & Discard item $e$ immediately & $[a, d, c]$ \\
    \addlinespace
    $6$ & $f$ & $k = 1 \in \{1 \dots 6\}$ & True ($1 \le 3$) & Replace slot $R[1] \leftarrow f$ & $[f, d, c]$ \\
    \addlinespace
    $7$ & $g$ & $k = 4 \in \{1 \dots 7\}$ & False ($4 > 3$) & Discard item $g$ immediately & $[f, d, c]$ \\
    \bottomrule
  \end{tabular}%
  }
\end{table}
\end{example}

\FloatBarrier
\section{Formal Proof by Induction of Reservoir Sampling}
\label{sec:reservoir-sampling-proof}

The mathematical correctness of Reservoir Sampling ensures that upon stream termination at any point, every observed element has an identical inclusion probability.

\begin{theorem}[Correctness of Reservoir Sampling]
For any stream $S$, at any step $n \ge m$, every element $i_j$ seen so far ($1 \le j \le n$) belongs to the reservoir $R$ with probability:
\begin{equation}
  \mathbb{P}(i_j \in R) = \frac{m}{n}.
\end{equation}
\end{theorem}

\begin{proof}
We proceed by mathematical induction on observed stream length $n$, keeping capacity $m$ fixed.

\textbf{Base step ($n = m$):}\\
After the first $m$ steps, elements $i_1, \dots, i_m$ are deterministically placed into reservoir $R$:
\begin{equation}
  R = [i_1, i_2, \dots, i_m].
\end{equation}
Therefore, for every $j \in \{1, \dots, m\}$:
\begin{equation}
  \mathbb{P}(i_j \in R) = 1 = \frac{m}{m}.
\end{equation}
The base case holds.

\textbf{Inductive hypothesis:}\\
Assume the property holds at step $n - 1$, so for every $i_j$ with $1 \le j \le n - 1$:
\begin{equation}
  \mathbb{P}(i_j \in R \text{ at step } n - 1) = \frac{m}{n - 1}.
  \label{eq:induction-hypothesis}
\end{equation}

\textbf{Inductive step (transition from $n - 1$ to $n$):}\\
Upon arrival of the $n$-th element $i_n$, the algorithm draws:
\begin{equation}
  k \sim \operatorname{rand}(1, n).
\end{equation}
We evaluate two mutually exclusive cases:
\begin{itemize}
  \item \textbf{Case 1: Newly arrived element ($j = n$):}\\
  Element $i_n$ enters reservoir $R$ if and only if $k \le m$. Because $k$ is drawn uniformly from $\{1, 2, \dots, n\}$:
  \begin{equation}
    \mathbb{P}(i_n \in R) = \mathbb{P}(k \le m) = \frac{m}{n}.
  \end{equation}
  The assertion holds for $j = n$.

  \item \textbf{Case 2: Any earlier element ($j < n$):}\\
  Consider an arbitrary previous element $i_j$ ($1 \le j \le n - 1$). For $i_j$ to remain in $R$ after step $n$, two conditions must hold:
  \begin{enumerate}
    \item $i_j$ was present in reservoir $R$ at step $n - 1$;
    \item $i_j$ \textbf{survives} step $n$, meaning it is not overwritten by $i_n$.
  \end{enumerate}
  Applying the conditional probability law:
  \begin{align}
    \mathbb{P}(i_j \in R \text{ at step } n) &= \mathbb{P}(i_j \in R \text{ at step } n - 1) \notag\\
    &\quad \cdot \mathbb{P}(i_j \text{ survives step } n \mid i_j \in R \text{ at step } n - 1).
  \end{align}
  By the inductive hypothesis~\eqref{eq:induction-hypothesis}, the first factor is $\frac{m}{n - 1}$.

  We evaluate survival event $E_{\text{surv}}$, meaning $i_j$ is not overwritten at step $n$. An element in $R$ survives under two disjoint scenarios:
  \begin{itemize}
    \item \emph{Scenario A:} The new element $i_n$ does not enter $R$ ($k > m$), occurring with probability:
    \begin{equation}
      \mathbb{P}(k > m) = 1 - \frac{m}{n}.
    \end{equation}
    Here, no slot is modified, and $i_j$ survives with certainty.
    \item \emph{Scenario B:} The new element $i_n$ enters $R$ ($k \le m$, probability $\frac{m}{n}$), \textbf{BUT} the replacement slot chosen is distinct from the one occupied by $i_j$. Because slot $k$ is chosen uniformly among the $m$ slots, the probability of selecting one of the other $m - 1$ slots is:
    \begin{equation}
      \frac{m - 1}{m}.
    \end{equation}
  \end{itemize}
  Summing probabilities of these disjoint outcomes gives total survival probability:
  \begin{align}
    \mathbb{P}(E_{\text{surv}}) &= \left(1 - \frac{m}{n}\right) + \left(\frac{m}{n} \cdot \frac{m - 1}{m}\right) \notag\\[1ex]
    &= \frac{n - m}{n} + \frac{m - 1}{n} \notag\\[1ex]
    &= \frac{n - m + m - 1}{n} \notag\\[1ex]
    &= \frac{n - 1}{n}.
    \label{eq:survival-probability}
  \end{align}
  Combining the survival probability~\eqref{eq:survival-probability} with the inductive hypothesis:
  \begin{align}
    \mathbb{P}(i_j \in R \text{ at step } n) &= \left(\frac{m}{n - 1}\right) \cdot \left(\frac{n - 1}{n}\right) \notag\\[1ex]
    &= \frac{m}{n}.
  \end{align}
  The factor $(n - 1)$ cancels algebraically, confirming the result for all prior items $j < n$.
\end{itemize}
Having established the assertion for both $j = n$ and all $j < n$, the inductive step is complete:
\begin{equation}
  \mathbb{P}(i_j \in R) = \frac{m}{n}, \qquad \forall j \in \{1, 2, \dots, n\}.
\end{equation}
The correctness of Reservoir Sampling is established.
\end{proof}

\FloatBarrier
\FloatBarrier
\section{Comparative Synthesis}
\label{sec:comparative-synthesis}

To consolidate our theoretical findings, \autoref{tab:sampling-comparison-streaming} summarizes the operational profiles, working space, and computational costs of the three stream sampling paradigms examined.

\begin{table}[H]
  \centering
  \caption{Comparative synthesis of stream sampling strategies.}
  \label{tab:sampling-comparison-streaming}
  \resizebox{\textwidth}{!}{%
  \begin{tabular}{l c c c c l}
    \toprule
    \textbf{Algorithm} & \textbf{Knowledge of $n$} & \textbf{Working Space} & \textbf{Cost per Item} & \textbf{RNG Calls} & \textbf{Data Structure} \\
    \midrule
    Algorithm S (Knuth) & Known a priori & $\mathcal{O}(1)$ & $\mathcal{O}(1)$ & $n$ ($1$ per item) & Static array for sample \\
    \addlinespace
    Priority Sampling & Unknown & $\mathcal{O}(m)$ & $\mathcal{O}(\log m)$ & $n$ ($1$ per item) & Bounded Min-Heap of size $m$ \\
    \addlinespace
    Reservoir Sampling (Alg. R) & Unknown & $\mathcal{O}(m)$ & $\mathcal{O}(1)$ & $n - m$ ($0$ for first $m$) & Reservoir array of size $m$ \\
    \bottomrule
  \end{tabular}%
  }
\end{table}

\paragraph{Algorithm Engineering Remark (Optional Insight: Practical RNG Bottleneck):}
In high-throughput industrial streaming workloads ($n \gg 10^9$), classical Algorithm R incurs a subtle computational penalty: although operating in $\mathcal{O}(1)$ worst-case time per item, it queries the pseudo-random number generator (RNG) for each of the $n - m$ subsequent stream items. Because admission probability $\frac{m}{j}$ rapidly approaches zero, nearly all RNG invocations yield $k > m$, wasting substantial CPU cycles on discarded items.
To circumvent this bottleneck, \textbf{Vitter's Algorithm Z}~\cite{vitter1985} models the number of elements skipped before the next reservoir replacement as a random variable termed the \emph{skip distance}. By sampling this skip distance directly, the algorithm skips entire blocks of unselected stream elements in bulk without per-item RNG queries, drastically cutting total random draws from $\mathcal{O}(n)$ to $\mathcal{O}(m(1 + \log(n/m)))$.
